\documentclass[../main.tex]{subfiles}
\begin{document}
%==========================================TIÊU ĐỀ TẬP========================================
\begin{table}[h]
    \begin{adjustwidth}{-1cm}{}
        \begin{tabular}{>{\centering\arraybackslash}p{6cm}|>{\raggedright\arraybackslash}p{12.5cm}}
            \multirow{5}{6cm}
            % Cột bên trái: ÉPISODE và số tập
            {\rainbowcolor
                {\\[-40pt]
                    \flaregothic\fontsize{20pt}{20pt}\selectfont \centering
                    \color{eptype}JOUR\color{black}\\[12pt]
                    \hbrk\fontsize{72pt}{72pt}\selectfont
                    \color{epnum}3
                    \\[18pt]
                    % \flaregothic\fontsize{14pt}{14pt}\selectfont
                    % \color{epattr}BIENVENUE, \uline{\color{Banpire} Mel} !
                    % \flaregothic\fontsize{18pt}{18pt}\selectfont
                    % \color{epattr}ÉPISODE PERDU\\
                    \unrainbowcolor}
                \color{black}}
             &
            % Dòng thứ nhất: tên tập tiếng Nhật
            {\fotskip\fontsize{18pt}{18pt}\selectfont にじさんじ\ruby{危}{き}\ruby{機}{き}\ruby{一}{いっ}\ruby{髪}{ぱつ}
            }                        \\[6pt]
            % Dòng thứ hai: tên tập tiếng Anh
             & {\cafeteria\fontsize{18pt}{18pt}\selectfont An Explosive Incident}                                               \\[4pt]
            % Dòng thứ ba: tên tập tiếng Việt
             & {\vnmsans\fontsize{18pt}{18pt}\selectfont Con tin}                                                               \\[8pt]
            % Dòng thứ tư: Nhân vật xuất hiện
             & {\sen{\fontsize{14pt}{14pt}\selectfont Track used: Shock Temperament (Faster Remix, from Professional Block 2)}}
            \\[6pt]
            % Dòng cuối cùng: Thời gian diễn ra của tập (thường sẽ là ngày phát hành)
             & {\continuum\fontsize{16pt}{16pt}\selectfont Ngày phát hành: 22/05/2021}                                          \\[18pt]
        \end{tabular}
    \end{adjustwidth}
\end{table}
%=========================================TRUYỆN CHÍNH========================================
\color{black}

\pov{Hikawa Kaigou}

Sáng ngày thứ ba ở thành phố Cầu Vồng.

Hôm nay trời lại đẹp. Không một gợn mây, nắng không quá gắt, gió xuân nhẹ như ve vuốt. Mọi thứ thật yên bình\dots\ nếu không tính đến cái cột nước khổng lồ hôm qua mà tôi giả vờ không biết gì.

``\eng{Kaigou-user, tôi thấy cô hồi hộp hơn thường ngày khi bước vào cửa hàng tiện lợi đấy. Có điêu gì khiến cô lo lắng sao?}'' giọng Game - cái thiết bị đeo tay do Kuon-kun đưa - vang lên đều đều, pha lẫn chút trêu chọc.

Tôi nheo mắt nhìn tấm kính của cửa, phản chiếu lại hình ảnh một thiếu nữ cao gầy, tóc đen xù buộc lệch, mặc áo khoác đen và quần ngắn sọc - những bộ mà tôi mượn được từ công ty. Trông như một cô chị khó gần, mà thật ra thì\dots\ tôi đúng là thế thật.

``\eng{Không có gì. Chỉ là\dots\ tôi không thích mấy cái chỗ đông người. Đặc biệt là nơi nào có nhiều đực rựa,}'' tôi đáp, lạnh tanh.

``\eng{Vẫn cái thái độ cũ nhỉ? Cô biết Kuon-user từng bảo gì không? `Kaigou-chan mà thấy trai là như thấy gián trong bếp!'} Game tiếp lời, giọng ngập tràn hài lòng tới mức khó hiểu.

Tôi rít qua kẽ răng: ``\eng{Đừng có mà nhắc tên cậu ta với mấy thằng đó trong cùng một câu.}''

``\eng{Tôi không nhắc tên `mấy thằng đó', tôi chỉ-}''

``\eng{Câm mồm. Một ngày không khịa tôi là không chịu được hả?}''

``\eng{Ờ thì\dots\ nếu cô gọi đây là tương tác xã hội, tôi cũng không phản đối.}''

Tôi thở dài. Từ lúc tôi quyết định đi cùng Kuon-kun sang thế giới này, tôi đã tự nhủ sẽ hạn chế tính khí quá khích. Nhưng mà\dots\ thật sự khó chịu. Mấy tên đàn ông ở cái công ty ANYCOLOR này trông nhếch nhác đến mức khiến tôi phát sợ. Nhất là cái ông tóc bạc Maimoto gì đó. Gặp hôm qua, ông ta cười cái kiểu như thể tôi là mớ hành lá cuối cùng còn sót lại trong chợ.

Tôi khựng lại khi Game đột nhiên lên tiếng: ``\eng{Cô biết đấy, họ đã liên hệ với Kuon-user để điều tra vụ hôm qua rồi.}''

``\eng{\dots Biết.}'' Tôi đáp, khẽ siết túi giấy trong tay. ``\eng{Nhưng tôi không nói gì cả. Vì an toàn của cậu ta và Suzu.}''

``\eng{Giá như cô nói năng dịu dàng với tôi được như với cậu ấy thì đời tôi đỡ cực nhỉ\dots}'' ông ta lầm bầm.

``\eng{Ông thì khác. Kuon-kun là người duy nhất trong cái thế giới chó chết này mà tôi còn thấy tin được.}''

``\eng{Ơ, tôi tưởng tôi cũng là-}''

``\eng{Ông thì là cái đồng hồ biết nói.}''

``\eng{Biết rồi, không cần nhấn mạnh vậy đâu\dots}''

Tôi sắp băng qua đường thì một giọng nữ quen quen vang lên từ xa, có phần hốt hoảng. ``\textbf{N-Nguy to rồi!}''

\begin{wrapfigure}[24]{r}{7cm}
    \includegraphics[width=7cm]{../../../../Image/Book?/chihiro_model.png}
\end{wrapfigure}

Một cô bé trông như học sinh tiểu học, nhưng tôi biết rất rõ đó là một cô gái mười tám tuổi mà tôi vẫn không biết vì sao.

\chrint

Yuuki Chihiro (\ruby{勇}{ゆう}\ruby{気}{き}ちひろ, \ruby{Du}{Dũng}-\ruby{ki}{Khí}\ Chi-hi-rô). Theo những gì mà cô bé giới thiệu, cô là một đứa trẻ 10 tuổi học tiểu học và là một nữ pháp sư, thứ mà tôi thấy ở các bộ hoạt hình trên TV. Cô bé là một trong những người đầu tiên gia nhập vào đại gia đình nhà Cầu Vồng này.

Chihiro-chan là một trường hợp khá đặc biệt trong cái thế giới này. Dù nhỏ tuổi, cô ấy lại có sức mạnh ma thuật đủ để khiến mọi người phải dè chừng. Đôi khi tôi thấy cô bé hóa lớn, thậm chí biến thành đàn ông chỉ để trêu mọi người. Đúng là không thể đoán trước được.

Nhìn bề ngoài thì cô bé lúc nào cũng vui vẻ, ngọt ngào, nhưng thực ra cô ấy khá hướng nội. Cô ấy luôn cố gắng kết nối với khán giả, kể cả khi bất đồng ngôn ngữ. Tôi từng thấy cô bé trả lời fan bằng tiếng Anh, tiếng Trung, dù không thực sự hiểu hết. Có lẽ vì thế mà cô ấy được yêu quý đến vậy.

Chihiro-chan nổi tiếng với những buổi stream kéo dài hàng chục tiếng, đặc biệt là chơi Apex Legends. Có lần tôi nghe nói cô ấy stream liên tục hơn 12 tiếng, vừa chơi vừa trò chuyện, có lúc còn khóc vì thất vọng với bản thân. Nhưng rồi lại cười, lại động viên mọi người. Đúng kiểu người không bao giờ bỏ cuộc.

Sau này, tôi nghe nói cô ấy đã rời Nijisanji vào ngày 31 tháng Một năm 2024. Dù tiếc nuối, nhưng tôi nghĩ cô ấy sẽ vẫn tiếp tục theo đuổi những điều mình yêu thích, dù ở bất cứ đâu. Với tính cách ấy, chắc chắn Chihiro-chan sẽ không bao giờ biến mất khỏi thế giới này.

\begin{wrapfigure}[35]{l}{6.5cm}
    \includegraphics[width=6.5cm]{../../../../Image/Book?/keisuke_model.png}
\end{wrapfigure}

Chihiro-chan có mái tóc dài màu dừa cạn (khoảng giữa màu xanh lam và tím) được tết thành hai bím cao bằng ruy-băng lớn màu dừa cạn, nhưng nghiêng nhiều về tím hơn, trên đỉnh có chỏm tóc ngố với đôi mắt xanh. Cô ấy mặc một chiếc váy có hai màu chủ đạo là trắng và dừa cạn với ruy băng và diềm xếp nếp, đi giày màu dừa cạn trông như có quai và tất trắng cao tới quá đầu gối. Phụ kiện mà Chihiro-chan đi kèm là đeo găng tay trắng với đường viền cổ tay màu hồng.

\chrint

``\textbf{Ở văn phòng chúng ta có biến rồi!}'' Chihiro-chan hét lớn. ``Ta đã bị bao vây!''

Tôi nhướng mày: ``Bao vây?'' Tôi quay sang Game, chỉ thấy màn hình ông ta nháy chữ ``Đang xử lý dữ liệu\dots''

Một giọng nam vang lên sau tôi, chính là người mà tôi nhắc tới ban nãy.

\chrint

Maimoto Keisuke (\ruby{舞}{まい}\ruby{元}{もと}\ruby{啓}{けい}\ruby{介}{すけ}, \ruby{Ma-i}{Vũ}-\ruby{mô-tô}{Nguyên}\ \ruby{Kê-i}{Khởi}-\ruby{xư-ke}{Giới}). Một ông chú cựu nông dân 36 tuổi không vợ không con và có tóc bạc trước tuổi.

Tôi không hiểu nổi tại sao cái công ty này lại nhận một ông chú như Maimoto Keisuke vào làm. Ông ta từng là nông dân, nhưng giờ chỉ biết ăn không ngồi rỗi, suốt ngày khoe khoang mấy chuyện tình cảm thất bại. Đàn ông kiểu này đúng là chẳng có gì ngoài cái miệng. Nghe nói ông ta từng có cả chục bạn gái, nhưng chẳng ai chịu nổi cái tính lười nhác và vô duyên của ông ta lâu hơn vài tháng, và rồi kết quả đều chỉ quy về một mẫu số chung: thất bại thảm hại. Đúng là mẫu đàn ông mà tôi chỉ muốn tránh xa.

Ông ta hay xuất hiện trong các buổi trò chuyện, lúc nào cũng cố tỏ ra hài hước, nhưng thực ra chỉ toàn mấy câu đùa nhạt nhẽo mà mấy ông già hay nói. Có lần tôi nghe Kuon-kun bảo ông ta từng hợp tác với vài người bên ngoài Nijisanji, thậm chí còn nghe cậu ta nói là có quan hệ (không phải kiểu tình cảm) với Suabru-chan bên \hll\ qua Shigure Ui. Chậc. Đàn ông mà không có chí tiến thủ, chỉ biết dựa dẫm vào người khác thì sớm muộn cũng bị bỏ lại phía sau thôi.

Tôi còn nghe nói ông ta hay đóng vai ``ông chú'' trong các chương trình, lúc nào cũng thích làm trung tâm chú ý, nhưng thực ra chẳng ai thèm quan tâm, thậm chí còn bị đem ra làm trò tiêu khiển. Đúng kiểu người sống dựa vào quá khứ, không chịu nhìn về phía trước. Nếu không phải vì công việc, tôi chẳng bao giờ muốn dây dưa với loại người như thế này.

Trong con mắt tôi, ông ta chẳng có một cái  gì để tôi có thể khen được cả.

Trong bộ trang phục thường ngày, Keisuke luôn mặc một bộ quần áo Nhật Bản truyền thống, với bộ áo kimino xanh hở ngực vạm vỡ (tôi có thể dám chắc đó không phải là giả), với quần xanh và dép obi. Ông ta có mái tóc bạc trắng và có đôi mắt xám, thỉnh thoảng tôi còn thấy ông ta đeo kính nâu nữa.

\chrint

``Ể!? Văn phòng bị bao vây!? Có địch à!?''

``Không phải địch đâu! Mà\dots\ còn nữa!'' Chihiro thở hồng hộc, ``Chaika đã bị bắt làm con tin rồi!''

Tôi mở lớn mắt. ``Ý em là cái gã tai nhọn lông mày rậm hôm bữa chị với Kuon-kun gặp ở quán ăn gia đình hôm kia hả?'' Thật, tôi vẫn không thể nào thẩm được cái gu ăn mặc của gã ta.

``Anh ấy đó!'' cô bé gật đầu.

Game lên tiếng, giọng giờ đã nghiêm túc: ``Giờ thì nghe có vẻ nghiêm trọng thật.''

``Rốt cuộc ai bao vây văn phòng vậy?'' tôi hỏi.

Chihiro không đáp. ``\textbf{Không còn thời gian đâu! Phải cứu anh ấy ngay!}''

Tôi khẽ cau mày, thở ra một tiếng.

``Dù gì thì ảnh cũng là người đã chỉ chỗ ở tạm cho mình, Suzu và Kuon-kun\dots'' tôi lẩm bẩm, rồi nhìn Game, nghiêm túc hơn.

``Chuẩn bị dẫn đường đi. Nếu mấy tên đó mà chạm vào bất cứ ai\dots\ thì tôi sẽ xé xác chúng.''

\begin{center}
    \uline{Văn phòng - Nijisanji.}
\end{center}

Văn phòng Nijisanji hôm đó chẳng khác gì một công trường đang thi công dở dang. Khắp nơi là hộp giấy chưa kịp dọn, đồ đạc ngổn ngang như bãi chiến trường. Chả ai buồn lau sàn nữa thì phải, vì tôi vừa bước vào đã thấy đế giày dính bụi bặm. Những chiếc hộp các-tông còn ngổn ngang khắp nơi, một số vừa được khui, giấy gói rơi vãi đầy sàn. Trên kệ tạm kê sát tường là những mô hình nhân vật, phụ kiện biểu diễn, bảng tên, một vài bộ quần áo trình diễn còn treo dang dở trên mắc, và vài tủ tài liệu cùng mấy chiếc bàn máy tính đặt sát vách. Căn phòng này rõ ràng là chưa chuẩn bị xong để đưa vào sử dụng.

Ở giữa phòng, một cô nàng có mái tóc hai màu bạc và đen, chia theo hai nửa rõ rệt, đang giữ lấy một cậu yêu tinh cao lớn - người hiện giờ đang hoàn toàn bất động. Tay còn lại của cô nắm chặt chiếc công tắc nút bấm với gương mặt đầy sát khí.

\chrint

\begin{wrapfigure}[22]{r}{11cm}
    \includegraphics[width=11cm]{../../../../Image/Book?/rena_model.png}
\end{wrapfigure}

Cô gái có tóc hai màu mỗi nửa trông như thể hóa trang bị lỗi kia là Yorumi Rena (\ruby{夜}{よる}\ruby{見}{み}れな, \ruby{Giô-rư}{Dạ}-\ruby{mi}{Kiến}\ Rê-na), một người gia nhập Nijisanji với cậu chàng CEO hôm trước kia. Cô là một nữ pháp sư kiêm idol 17 tuổi đã tạ thế tới thế giới này để thu hút những người hâm mộ của mình ở đây.

Yorumi Rena là một trường hợp khá đặc biệt trong Nijisanji. Cô ấy nổi tiếng với tính cách ``khó đoán'' và đôi khi hơi ``quái dị'' đến mức khiến người khác phải dè chừng. Tôi từng nghe mấy người trong công ty gọi cô là ``Yorumi Rena Psychopath'' nhưng không phải theo kiểu bệnh lý, mà là kiểu người có thể chuyển trạng thái cảm xúc cực nhanh, lúc thì vui vẻ, lúc lại lạnh lùng, thậm chí có thể nổi điên bất chợt. Có lần tôi thấy cô ấy vừa cười nói với mọi người, quay đi một cái đã lẩm bẩm mấy câu nghe rợn cả gáy.

Cô ấy là một kiểu người rất nhát ma. Theo lời của Hayato, có lần con nhỏ phải kêu cứu cậu ta trong một buổi chơi RE:7 vì quá sợ hãi. Dường như kinh dị không phải là sở trường của cô ấy.

Tuy vậy, Rena-chan lại rất quan tâm tới bạn bè, nhất là những người mới vào công ty. Cô ấy thường chủ động giúp đỡ, hướng dẫn mọi người, nhưng cách thể hiện thì chẳng giống ai, kiểu vừa giúp vừa trêu, khiến người ta không biết nên cảm ơn hay nên tránh xa. Có lần tôi nghe nói cô ấy từng tổ chức một buổi stream chỉ để ``bắt cóc'' một đồng nghiệp rồi tra khảo như phim trinh thám, cuối cùng lại tặng quà cho người đó.

Nói chung, Yorumi Rena là kiểu người mà nếu không cẩn thận sẽ bị cuốn vào mấy trò quái gở của nhỏ lúc nào không hay. Nhưng nếu hiểu được tính cách ấy, thì lại thấy cô ta rất thú vị và đáng quý. Dù có hơi ``tâm thần'' theo nghĩa vui, nhưng Rena-chan luôn biết cách làm cho mọi người xung quanh cảm thấy đặc biệt, dù là bằng cách nào đi nữa.

Rena-chan có mái tóc buộc hai bím, với bên trái là màu đen và bên phải màu trắng. được buộc cao bằng ruy băng cánh hoa đen lớn (nhưng lần này tôi thấy cái ruy băng cánh hoa kia là màu cầu vồng lòe loẹt), cùng có vài sợi tóc được nhuộm hồng và buộc ở đuôi tóc bằng ruy băng. Cô nàng có đôi mắt đỏ.

\begin{wrapfigure}[32]{l}{7cm}
    \includegraphics[width=7cm]{../../../../Image/Book?/chaika_model.png}
\end{wrapfigure}

Cô nàng mặc một chiếc váy xếp nếp màu trắng với tay áo ngắn bồng bềnh, một chiếc váy lớp ngoài mang hai tông màu trắng-đen được buộc bằng ruy băng ở giữa bụng, đeo găng tay trắng dài đến khuỷu tay và quần tất dài đến đùi gồm màu trắng với viền thường bên phải và màu đen có họa tiết thỏ bên trái. Cô ấy đi đôi mary-janes với dây quai đan chéo, chân phải có quấn thêm một sợi ruy băng đen nữa.

Người còn lại đang nằm trên tay của cô ấy và đang bất tỉnh kia là Hanabatake Chaika (\ruby{花}{はな}\ruby{畑}{ばたけ}チャイカ, \ruby{Ha-na}{Hoa}-\ruby{ba-ta-kê}{Đèn}\ Chai-ca). Gã này đã từng sở hữu một quán bar ở thế giới trước, nhưng không hiểu vì sao mà lại lạc lõng tới đây.

Và bây giờ, ở thế giới Nijisanji này, gã sở hữu một quán cà phê, cách chúng tôi sống không quá xa. Theo như những lời giới thiệu ở công ty, gã này có một trực giác nhạy bén và rất thông minh, vậy mà giờ đây, lại bị một nữ pháp sư ỏng ẹo kia bắt được.

Nghe nói gã ta là ``thủ lĩnh'' của cái gọi là ``Binh đoàn kháng chiến Nijisanji'' ngoại trừ việc gã không đánh đấm gì cả, mà chỉ là cái cớ để lập nhóm. Nói muốn làm bạn thì nói luôn đi, cứ khúm núm làm gì không biết?

Tôi còn nghe người ta đồn rằng gã đã từng ``tự sướng'' bằng nước keo vì muốn ``quan hệ'' với các cô gái tới mức mãnh liệt, thậm chí còn\dots\ coi một cái cuộn tóc làm ``lỗ hoa'' nữa. Thảm hại không buồn nói.

Được cái là gã ta đã giới thiệu cho chúng tôi chỗ để chúng tôi có thể nghỉ ngơi - cũng là nơi mà tôi, Suzu và Kuon-kun ở đó - chỉ để rồi gã ấy giở giọng muốn làm quen gần gũi với tôi. Và bạn biết kết cục sau đó là gì rồi đấy: gã bị tôi hạ nốc-ao tại chỗ, tới mức phải để Chihiro-chan và lão Keisuke kia khiêng về.

Chaika có mái tóc đen với mái nhuộm xanh lá, có vài cọng được dựng đứng lên bởi keo. Gã ta có đôi mắt vàng, trang điểm đậm và tô son xanh lá, đeo khuyên tai như phụ nữ. Thậm chí bộ trang phục của gã mặc cũng là của phụ nữ: bộ hầu gái mang hai màu đen-trắng ngắn tay, có viên ngọc xanh lá đính ở ngực áo, đeo găng tay với vòng đen có lông, đi quần tất và đi giày đen có đính ngọc ở phía trước. Cái kiểu ăn mặc nam không ra nam mà nữ không ra nữ như thế này khiến tôi khó chịu.

\chrint

\img{nj-3a}

``\textbf{Ở trong đây nè!}'' Chihiro-chan kêu to, chỉ tay vào phòng trong. Mọi người gần như đồng loạt xông vào - không ai che giấu nổi vẻ lo lắng.

``Chaika! Cậu ổn chứ!? Bọn tôi đến cứu cậu đây!'' lão Keisuke la lên như phim hành động hạng B.

``Có ai giải thích cho tôi chuyện-'' tôi bắt đầu hỏi, nhưng ngay khi mắt cô chạm vào Chaika đang bị kẹp dưới tay Rena-chan, miệng đang run run như sắp chết đói, thì cô bỗng khựng lại.

Tôi khoanh tay lại, chưa kịp hỏi chuyện gì đang diễn ra thì cảnh trước mắt đã đủ cho tôi hiểu vấn đề. ``\dots Cái éo gì đây?'' tôi thốt ra trong tiếng nghiến răng.

``Err\dots\ Yorumi-san, cô đang làm cái trò gì vậy?'' Game lên tiếng, vẫn cái giọng dửng dưng đó.

``\textbf{Chị vừa làm gì anh ấy vậy!?}'' Chihiro-chan bên cạnh gào lên.

``\textbf{Đừng có lại gần đây!}'' Rena-chan hét lên, giơ cái công tắc run rẩy như người mới tập múa võ. ``Nếu không thì tôi cho nổ banh cái phòng này đấy!''

Tôi đảo mắt, chống hông đáp lại: ``Ờ rồi, đe dọa đi, còn thiếu mỗi khăn trùm đầu và ghi âm đòi tiền chuộc nữa là đủ combo rồi đấy.''

Game vẫn không nhịn được: ``Trông cô còn run hơn cả thạch. Cầm cái nút bấm mà như cầm kem đang tan chảy thế kia.''

``Im cái miệng lại,'' Tôi liếc hắn. ``Không thấy chỗ này đầy thuốc nổ à? Ông mà khích con nhỏ đó làm liều thì liệu hồn.''

``Cô cũng thế còn nói ai\dots'' ông ta biện hộ khi gã Keisuke và Chihiro-chan còn đang lườm chúng tôi.

Chúng tôi lùi dần về đằng sau, để cho hai người kia cùng Cầu Vồng nói chuyện. Gã cựu nông dân cố đàm phán: ``Yorumi, tại sao cô lại làm chuyện này?''

\img{nj-3b}

``''\textbf{Mọi người không thấy gì à!?'}' cô nàng Ma pháp hét lớn, rồi ném cái công tắc thẳng vào người Chaika, bấp bênh giữa hai mặt, rồi rơi đúng ngay giữa ngực. Tôi phải công nhận: cú ném ấy mà trượt thì cũng có thể phát nổ thật.

``Nhìn cái phụ kiện tóc của em đây này!''

Tôi nheo mắt nhìn lên đầu nhỏ đó. Những chiếc ruy-băng màu sắc lòe loẹt, cùng với những cái nơ cài tóc, chúng đang hoàn toàn đổi màu theo thời gian. ``Ờm\dots\ nhìn cũng\dots\ khác thường đấy,'' tôi chưa muốn gây chuyện nên cố nhịn.

\img{nj-3c}

``Trông cũng có vẻ xịn xò mà,'' cô bé mười tuổi ngạc nhiên, khẽ trố măt.

``Xịn xò đâu mà xịn xò!??'' Rena-chan tức tối. ``Em nhờ Suzu làm cái phụ kiện dễ thương. Và em ấy đưa cho em cái thứ\dots\ này đây!''

Tôi hơi khựng lại. ``Suzu làm cho em sao?''

``Đúng vậy! Em bảo em ấy là làm cho dễ thương vào! Ai ngờ em ấy làm ra cái đồ của bà u mê chơi búp bê hồi năm 2000!''

Tôi khẽ thở dài. Suzu vốn khéo tay mà, nhưng gu nghệ thuật dôi khi cũng hơi\dots\ lạ, tôi phải thừa nhận. Dù sao thì cũng còn đỡ hơn mấy thằng con trai làm gì cũng vụng như mèo mù leo cây.

Dù vậy, nghe nhỏ này gào lên chê, tôi cũng hơi chột dạ. Không thể để nhỏ ấy xúc phạm em gái tôi ngay giữa phòng như thế này được.

Game chen ngang: ``Có thể Suzu-user đang thử nghiệm gì đó? Hay đơn giản là\dots\ cô không hiểu nổi nghệ thuật?''

``Ông khen được cái đó thì đúng là gu của ông cũng chuẩn hãm rồi đấy,'' tôi nhíu mày. ``Còn cái tên kia nữa\dots'' tôi liếc qua Chaika vẫn nằm xụi lơ.

Thề, chưa thấy ai mặc đồ hầu gái mà mặt thì rõ là mặt đàn ông. Trông cứ như cosplay rẻ tiền bị bỏ giữa đường. Hôm đầu mới gặp tôi còn suýt thốt ra: ``Cậu là gái à? Hay trai giả gái?''  nhưng vì có Suzu với Kuon-kun bên cạnh nên tôi ngậm miệng. Giữ thể diện cho người nhà là ưu tiên hàng đầu.

Kuon-kun lúc ấy chỉ cười trừ mà nói: ``Gu ăn mặc của cậu trông\dots\ ờm, đặc biệt quá.'' Nhẹ nhàng. Lịch sự. Còn tôi? Tôi mà nói thì chỉ có ``Trông như chọc mù con mắt'' là nhẹ nhất.

``Thế nên tôi bắt Chaika làm con tin!'' Rena lại hét lên. ``Cho đến khi Suzu làm lại một cái dễ thương hơn! Bằng không thì\dots''

``Cô lại dọa nổ phòng à?'' tôi lầm bầm. ``\textit{Bài đó xưa như Diễm rồi\dots}''

``Không.'' Nhỏ đó lấy ra\dots\ một xiên oden?

Tôi chớp mắt.

``Ờm\dots\ thế rồi sao?'' tôi nhíu mày.

``Tôi sẽ bón cái này cho Chaika!''

\img{nj-3d}

``\dots Hả???'' Cả tôi lẫn Game đều đứng hình. Còn Keisuke thì giật thót như thể vừa bị sét đánh ngang tai. ``Cô định\dots\ ép cậu ta ăn sao!?''

``Đúng vậy! Nếu không đáp ứng yêu cầu, tôi sẽ cho Chaika ăn một bữa đẫy ắp đồ nóng sốt!''

``\dots Thế rốt cuộc cái nút bấm kia để làm cái quái gì?'' Game nhìn sang Chaika.

``Đói\dots\ quá\dots'' cậu ta thì thào, mặt tái mét như sắp hóa zombie đến nơi. ``Có ai\dots\ cho tôi\dots\ ăn đi\dots''

Tôi quay đi, ngửa mặt thở dài và tặc lưỡi một tiếng. ``Cái thế giới này đúng là không thể hiểu nổi\dots''

``Cậu ta đói đến mức nói sảng rồi!'' Chihiro-chan cuống lên.

``Thế éo nào mà cậu ta lại yếu thế được nhỉ? Nhìn cao to thế kia mà\dots''

``Không đơn giản đâu,'' Game gãi đầu. ``Bình thường Chaika đâu như vậy. Có khi là\dots''

``Có khi cậu ta sắp bung nội tại biến hình thành quái vật ăn thịt người cũng nên,'' tôi lầm bầm.

``Không ai biết chắc điều gì sẽ xảy ra nếu cậu ta ăn một miếng oden cả,'' cô bé Ma pháp sợ hãi.

Tôi nhìn kỹ lại căn phòng, rồi liếc sang nút công tắc vừa bị quăng, thuốc nổ còn lăn lóc khắp sàn, Chaika nằm vật như cái bánh bao thiu, và con nhỏ Rena-chan thì giơ xiên đồ ăn như đang hành hình. Tôi vừa nhìn vừa không biết phải tức giận hay bật cười.

``Vậy thì còn chờ gì nữa!?'' cô ta gào lên. ``Mau gọi Suzu lại đây và bảo em ấy làm lại cho tôi phụ kiện tóc khác đi!''

Tôi lấy điện thoại ra, dù biết rõ cái viễn cảnh vô ích trước mặt. Cái máy của Kuon-kun thì tôi biết thừa: cậu ta vứt máy ở nhà vì quên sạc.

``Cô nghĩ cậu ta sẽ nghe máy sao?'' Game nói.

``Ể?'' Chihiro-chan nhướng mày.

``Tôi gọi rồi. Chỉ có tiếng tút tút. Máy ở nhà rồi.''

Tôi thở dài, bỏ điện thoại xuống. ``Sao cậu ta lại có thể quên sạc điện thoại được nhỉ\dots?''

``Vậy cô thử gọi cho cô em gái của cô xem?'' lão Keisuke thảng thốt.

``Con bé không có máy,'' tôi đáp gọn.

``V-Vậy thì\dots'' nhỏ Rena-chan dường như bắt đầu thương lượng. ``Kiếm cho tôi một cái phụ kiện khác cũng được! Nhưng nhất định nó phải dễ thương hơn cái của Suzu làm!''

Tôi đảo mắt. Lạy hồn, con nhỏ này đúng là dở hơi.

Trong khi đó, lão Keisuke thì bắt đầu cái màn đàm phán kiểu người lớn: ``Yorumi, nghe tôi nói này. Không cần phải làm những chuyện phức tạp như vậy đâu.''

Tôi khoanh tay, nhìn về phía cô nàng hai màu tóc  kia. ``Mà này, tôi hỏi thật, lúc em nhờ Suzu là em nói cụ thể thế nào vậy?''

``Thì\dots'' cô ta bĩu môi, ``Em bảo là `làm cho em một cái ruy băng sao cho nó trông đáng yêu là được!' Đại khái thế.''

``\dots Bảo sao.'' Tôi chép miệng.

``Cô nói kiểu đó thì đến tôi cũng chịu,'' Game chen vào. ``Nói kỹ hơn tí đi chứ!''

``Chính xác,'' lão cựu nông hùa theo. ``Nếu cô nói rõ ràng hơn thì Suzu chắc chắn sẽ làm tốt hơn đấy. Cứ mập mờ thế ai mà làm được?''

Lão nông dân kia tiếp tục ra sức cứu vãn tình hình, giọng điệu thì như đang kêu gọi hội cha mẹ học sinh: ``Chihiro, Chaika, Kaigou và tôi\dots\ tất cả chúng tôi sẽ cùng hỏi lại Suzu! Nên là\dots\ làm ơn đừng làm gì dại dột!''

Tôi và Chihiro-chan quay sang nhìn ông ấy với hai ánh mắt lạnh như kem. ``Khoan đã, cả tôi nữa á?'' tôi trừng mắt.

``\hl{Đừng để ý mấy cái tiểu tiết làm gì,}'' ông ta nói nhỏ như thỏ thẻ. ``Cứ giúp tôi thuyết phục Yorumi đã!''

\img{nj-3e}

``Nghe ông nói thế còn thấy mất thuyết phục hơn ấy,'' Game lắc đầu.

Cô bé nhỏ tuổi thì đành gật gù, miệng lí nhí: ``Đ-Được rồi\dots''

Tôi lắc đầu ngao ngán. ``Rồi, làm thì làm.''

Cả ba chúng tôi quay sang phía Rena, người mà giờ đây đang càng lúc càng bốc hỏa như ấm siêu tốc.

Keisuke lại bắt đầu tung chiêu cảm động kiểu phim truyền hình cũ rích: ``Yorumi, dừng lại đi! Cô không thực sự muốn làm vậy mà, đúng chứ?''

Tôi cũng gật đầu lấy lệ, nói thêm một câu theo bản năng: ``Cô mà không hạ cái xiên xuống là tôi cho cô ăn đấm đấy.''

``\dots Này, Kaigou-san?'' lão ta quay sang tôi với vẻ không tin nổi. ``Cô không định\dots''

``Đúng như ông nghĩ đấy,'' tôi đáp tỉnh bơ, quay cánh tay của mình. ``Tôi không nương tay đâu, kể cả là con gái.''

``Nhưng mà-'' Chihiro-chan định phản đối, nhưng Game thì chen vào, giọng chán đời: ``Trông mọi người lộn xộn quá\dots\ Tôi nghi là cô ta chẳng tin ai trong đám này đâu.''

Và rồi\dots\ rụt rịt.

Tiếng sụt sịt phát ra trong phòng. Tất cả mọi ánh nhìn đổ dồn vào con nhỏ Rena-chan.

Cô ta đang\dots\ khóc?

Đúng thế. Cô nàng tóc hai màu đó đang dụi mắt, môi mím chặt, còn tay thì vẫn cầm xiên oden. Nhìn vừa tội vừa ngớ ngẩn.

\img{nj-3f}

``\dots Tôi không nghĩ là cô ta thực sự cảm động vì mấy lời nhạt nhẽo đó,'' trợ lý trên đồng hồ của Kuon-kun lẩm bẩm, vẻ ngạc nhiên không giấu nổi trên giọng.

Tôi thì lặng người. ``\dots Không ngờ là cách của ông cũng có hiệu quả thật đấy, ông già.''

Keisuke nhẹ nhàng gọi: ``Yorumi à\dots'

``Yorumi\dots'' Chihiro-chan cũng gật đầu theo.

Rena-chan lau nước mắt, rồi nhìn ba người bọn họ với ánh mắt như sắp phát biểu một bài diễn văn giành giải Oscar: ``Nếu\dots\ nếu mọi người thực sự muốn giúp tôi như vậy\dots\ T-Tôi sẽ\dots''

Cô ta lại khóc tiếp.

\dots Ờ, tôi thì lại đang tự hỏi: tại sao bản nhạc ``Canon cung Rê trưởng\footnote{Tên gốc của bài là ``\textit{Bản luân khúc cung Rê trưởng cho ba đàn vĩ cầm và bè trầm đánh số}'', hay đơn giản là ``Canon in D'', là một trong những bản nhạc nổi tiếng nhất của Johann Pachelbel (1653 - 1706).}'' lại bắt đầu vang lên thế? Cái vẹo gì vậy? Ai mở nhạc? Ai đạo diễn cái khung cảnh đầy kịch tính giả tạo này?

Và rồi như thể số phận trêu ngươi, con nhỏ Rena-chan vô thức làm rơi cái xiên oden trong tay xuống trong khi đang lau nước mắt (tôi vẫn nghi là nước mắt cá sấu thì đúng hơn). Đồ ăn rơi xuống sàn là chuyện vặt, đúng không?

Không.

Cái que xiên kia rơi đúng chỗ\dots\ \emph{Cái. Nút. Bấm.}

Nó rơi kiểu chậm rãi, từ tốn như một đoạn tua chậm trong phim truyền hình Ấn Độ, và khi đầu gỗ trong tôi chưa kịp gào lên ``Đ-Đ-Ừ-N-G!!!'' thì\dots

*CẠCH*

Tiếng nút bấm vang lên như một cái búng tay huỷ diệt. Rồi sau đó là tiếng còi báo động gào rú như thể vừa mở tiệc khai tử cho cả toà nhà. Đèn nhấp nháy đỏ loá. Không khí thì náo loạn.

\img{nj-3g}

``\textbf{Á!}'' Rena-chan hét.

``\dots Cái nút bấm!'' Game kêu.

``Khoan đã, cái gì cơ!?'' tôi gào lên.

``Ể?'' Keisuke còn chưa hiểu mô tê.

``Chết rồi!'' Chihiro-chan tái mặt.

\begin{figure}[H]
    \centering
    \begin{subfigure}[t]{0.45\textwidth}
        \centering
        \includegraphics[width=8cm]{../../../../Image/Book?/nj-3h1.png}
    \end{subfigure}
    \quad
    \begin{subfigure}[t]{0.45\textwidth}
        \centering
        \includegraphics[width=8cm]{../../../../Image/Book?/nj-3h2.png}
    \end{subfigure}
    \quad
    \begin{subfigure}[t]{0.45\textwidth}
        \centering
        \includegraphics[width=8cm]{../../../../Image/Book?/nj-3h3.png}
    \end{subfigure}
    \caption{Phản ứng của mọi người khi chiếc nút tự hủy vô tình được kích hoạt.}
\end{figure}

\dots

\dots

*{\Large \textbf{BÙM!!!}}*

\img{nj-3i}

Cái văn phòng to uỳnh của Nijisanji nổ đùng một phát như vừa bị trúng tên lửa đạn đạo. Cột lửa bùng lên. Đất đá gầm rú như trâu điên, cả cái thành phố rung bần bật như đang lên cơn sốt. Từng mảnh kính, tường, sắt thép bắn tung tóe. Nếu có ai sống sót thì đúng là\dots\ kỳ tích.

Từ chỗ nào đó ngoài xa, tôi nghe thấy tiếng cô bé Ma thuật nhỏ tuổi thốt lên trên cao: ``Mấy cái kết nổ tung này đúng là tệ mà!''

Còn tôi thì chỉ kịp thở ra một tiếng: ``Chết dở thật rồi.''

Tệ hơn nữa là - cái vụ nổ đó không chỉ mình chúng tôi thấy. Thậm chí từ xa, Kuon-kun và Suzu cũng trông thấy cột khói bốc lên từ đống tro tàn.

``Tưởng chỉ mình thế giới \hll\ chơi trò nổ, ai ngờ ở đây cũng máu lửa ghê,'' cậu ta lẩm bẩm.

``Kiểu này là đền tiền không biết bao nhiêu mà kể,'' Game nói thêm như thể đang đọc hóa đơn điện nước.

Cả hai chúng tôi sau đó bị văng đi như những đám kia, nhưng nhờ cái mạch nước khốn kiếp ngày hôm qua mà ba cô gái đần độn không kém kia vô tình khơi ra, chúng tôi được một cột nước phun lên hất ngược trở lại không trung rồi rơi bịch xuống một đống thùng các-tông gần đó. Chấn thương nhẹ, ê mông một chút thôi, không chết.

Và đấy, khi tôi còn đang ho sặc sụa vì dính bụi tro, cái đầu tôi vẫn không quên quay qua chất vấn đồng đội của mình.

``\vie{Này, Suzu\dots\ cái phụ kiện chết tiệt hôm nay là do em làm hả?}''

``\vie{Ý chị là cái ruy-băng đổi màu `phê ngầu lòi' của Rena-nee á, Onee-sama?}'' Suzu vẫn còn dính bụi ở mặt do vụ nổ, nhưng miệng vẫn tỉnh bơ. ``Onii-sama làm cái đó đấy, dựa theo ý tưởng của chị ấy cơ.''

Kuon-kun cũng đứng kế bên, vẻ mặt chẳng lấy gì làm tự hào: ``\vie{Cô ấy cứ bảo là `em muốn có một cái phụ kiện tóc thật\~{} là đẹp', nhưng lại đếch thèm nói rõ như thế nào cả\dots}''

Game xen vào: ``\vie{Chuẩn chưa. Không biết rõ yêu cầu thì kết quả là thế đấy.}''

Tôi khoanh tay: ``\vie{Rồi, rồi. Thế là hai người các cậu bịa ra cái nơ đổi màu đó hả?}''

``\vie{Ừ thì,} cậu ta nhún vai. ``\vie{Tôi hỏi tới hỏi lui mà Rena-chan cứ khăng khăng phải `thật đáng yêu', `thật đặc biệt', `khiến cho người ta phải ngoái nhìn', không thêm gì khác. Thế là bực quá làm luôn ruy-băng đổi màu cho xong, còn đỡ bị nhây.}''

Suzu bật cười: ``\eng{Thực ra em cũng thấy cái nơ đó cu-te mà! Kiểu, mang hơi hướng hiện đại và ngầu ngầu, như trong game idol ấy!}''

Tôi thở dài. ``\vie{Tôi hiểu cảm giác của cậu mà, Kuon-kun\dots\ Mỗi lần làm việc với mấy người kiểu đó xong chỉ muốn đập đầu vào tường\dots}''

``\vie{Thế giờ tính sao?} Game hỏi.

Tôi ngó lại toà nhà sau lưng, giờ chỉ còn là đống tro tàn. ``\vie{\dots Tôi nghĩ đừng bao giờ để con nhỏ đó yêu cầu thêm phụ kiện nào nữa.}''
%==========================================TRUYỆN PHỤ=========================================
\omk

Tối hôm đó, trong cái không khí nửa âm u, nửa mệt rã sau một ngày dài như hàng thế kỷ, Kuon-kun đưa cho tôi một cái hộp nhỏ xíu, dài như hộp nữ trang. Mở ra, tôi thấy có một cái nơ buộc tóc cầu vồng, kiểu như dành riêng cho mấy nhỏ sặc sỡ hay đi hóa trang ngoài phố Harajuku.

Cũng chính là thứ mà nhỏ Rena-chan kia đeo trên đầu, nhưng được tinh chỉnh một chút để trông dễ nhìn hơn.

Tôi nhìn nó một lúc, rồi nhìn lại cậu con trai đang ngồi bó gối trong góc bàn, mắt liếc tôi chờ phản ứng như thể sắp bị ăn chửi đến nơi (tôi biết, tôi sẽ không làm thế với cậu ta, nhưng trông Kuon-kun khá hồi hộp). Tôi cầm nơ lên, xoay một vòng rồi buộc thử lên tóc.

``\vie{Hừm\dots}'' cậu ta gật đầu, ``\vie{của cô trông còn ổn áp hơn đấy.}''

``\vie{Cậu nói thật đấy à?}'' tôi nghiêng đầu, hơi nghi hoặc. ``\vie{Tôi tưởng cậu sắp chê nó giống giẻ lau cơ.}''

``\vie{Cũng hơi lòe loẹt với ố dề thật\dots}'' cậu ta nhún vai. ``\vie{Nhưng mà nó hợp với cô hơn con nhỏ Rena-chan nhiều.}''

Tôi cười chát. Không biết tôi có nên vui không nhỉ?

Từ trong gian bêp nhỏ tí kia, Suzu đang lúi húi dọn bát đĩa. Con bé nói nó muốn nấu món gì đó dễ ăn cho chúng tôi, sau cái ngày mà đám đồng nghiệp của chúng tôi hoá thành pháo bông trên trời. Thú thật, tôi không thấy đói mấy, nhưng nhìn đứa em tôi đang gấp rau củ bằng tay trần mà vẫn cười toe, tôi không nỡ từ chối.

Dù gì thì lâu lắm rồi tôi mới được ăn một bữa ăn ra trò do chính tay con bé chuẩn bị mà.

Tôi quay sang hỏi Kuon-kun: ``\vie{Bên ANYCOLOR có gọi không? Họ hỏi gì vụ cột nước hôm qua?}''

``\vie{Có,}'' cậu đáp, tay cầm cốc trà. ``\vie{Hỏi kiểu vòng vo thôi. Tôi chỉ bảo là hôm đấy ngủ dậy muộn, không biết gì hết. Chứ nói thật chắc bọn mình bị bắt rồi.}''

Tôi gật đầu, không ngạc nhiên. ``\vie{Biết thế nên tôi mới không đào sâu thêm. Bọn đó đúng là phiền thật.}''

``Ừ,'' tôi nói, nhìn lên cái đồng hồ treo tường cũ kỹ, ``\vie{tôi chỉ mong rời khỏi cái nơi này càng sớm càng tốt\dots}''

Suzu mang đồ ăn ra. Một món xào, một bát canh nóng, ba bát cơm trắng với một nồi cơm ú ụ đằng sau. Mùi thơm phảng phất gợi tôi nhớ về thời còn ở nhà.

Cả ba chúng tôi ngồi xuống, chắp tay: ``Itadakimasu.''

Tôi lặp lại lời ấy, khẽ cười với bữa cơm giản dị này. Kuon-kun thậm chí cũng không có ý kiến gì về món cơm do con bé làm, thậm chí còn ăn hẳn bốn bát cơm. Suzu thì khi thấy anh trai - dù là tự nhận - khen ngon, con bé cười tít mắt.

Miễn là cả ba chúng ta còn sống, vẫn còn ăn được với nhau, thì chắc vẫn chưa đến nỗi nào.

\pov{-}

Về số phận của những người khác\dots

Yuki Chihiro là người duy nhất trong số những ``nạn nhân vụ nổ'' vẫn còn xuất hiện sau đó. Cô được tìm thấy trong tình trạng bất tỉnh ở một khu dân cư nằm ở phía Bắc thành phố, nhưng không có thương vong gì quá nặng.

Còn lại, Rena, Keisuke và Chaika giống như Saku, Yuika, và Ririmu thì đều ``bay lên trời'' và từ đó biệt tăm biệt tích, không còn thấy đâu trong suốt những ngày tiếp theo. Không ai biết họ đã trôi dạt về đâu, chỉ có bầu trời phía Tây mỗi đêm lại sáng thêm một vì sao.

Thành phố nơi họ sống tạm giờ đây vắng đi một phần nhộn nhịp.

Và như thế, ngày thứ ba ở vùng đất Cầu Vồng khép lại trong hỗn loạn, mệt mỏi và tiếc nuối.

Một ngày mà tưởng chừng sẽ đơn giản chỉ là đi chợ, lại kết thúc bằng một vụ bắt cóc giả và vụ nổ thật. Một ngày mà con người bộc lộ đủ mọi cảm xúc: giận dữ, thương cảm, hoài nghi, rồi cuối cùng là chấp nhận.

Nhưng điều mà họ không biết là đây mới chỉ là khởi đầu.

Ngày mai, mọi thứ sẽ còn rối ren như vậy nữa.

\end{document}